% Shared layout of the experiment reports in reports/. Each report starts with
%   \documentclass[11pt,a4paper]{article}
%   % Shared layout of the experiment reports in reports/. Each report starts with
%   \documentclass[11pt,a4paper]{article}
%   % Shared layout of the experiment reports in reports/. Each report starts with
%   \documentclass[11pt,a4paper]{article}
%   % Shared layout of the experiment reports in reports/. Each report starts with
%   \documentclass[11pt,a4paper]{article}
%   \input{reports/style}
% (Overleaf compiles from the project root, so paths are relative to it.)
\usepackage[utf8]{inputenc}
\usepackage[T1]{fontenc}
\usepackage{lmodern}
\usepackage[margin=2.4cm]{geometry}
\usepackage{microtype}
\usepackage{amsmath,amssymb}
\usepackage{booktabs}
\usepackage{tabularx}
\usepackage{array}
\usepackage{enumitem}
\usepackage{graphicx}
\usepackage[table]{xcolor}
\usepackage{textcomp}
\usepackage{csquotes}
\usepackage{newunicodechar}
\newunicodechar{θ}{\ensuremath{\theta}}
\newunicodechar{φ}{\ensuremath{\varphi}}
\newunicodechar{λ}{\ensuremath{\lambda}}
\newunicodechar{Δ}{\ensuremath{\Delta}}
\newunicodechar{≈}{\ensuremath{\approx}}
\newunicodechar{°}{\ensuremath{^\circ}}
\newunicodechar{π}{\ensuremath{\pi}}
\newunicodechar{√}{\ensuremath{\surd}}
\newunicodechar{²}{\ensuremath{^2}}
\newunicodechar{³}{\ensuremath{^3}}
\newunicodechar{−}{\ensuremath{-}}
\newunicodechar{⋅}{\ensuremath{\cdot}}
\usepackage[most]{tcolorbox}
\usepackage[section]{placeins}
\usepackage{titlesec}
\usepackage{fancyhdr}
\usepackage{lastpage}
\usepackage[authoryear,round]{natbib}
\usepackage[colorlinks=true,linkcolor=blue!50!black,citecolor=blue!50!black,urlcolor=blue!50!black]{hyperref}

\setlist{itemsep=2pt,topsep=4pt}
\setlength{\parindent}{0pt}
\setlength{\parskip}{5pt}
\newcolumntype{L}{>{\raggedright\arraybackslash}X}
\newcolumntype{P}[1]{>{\raggedright\arraybackslash}p{#1}}

\definecolor{accent}{HTML}{2A78D6}
\definecolor{boxblue}{HTML}{EEF3FB}
\definecolor{trgrey}{HTML}{F4F4F1}
\definecolor{trline}{HTML}{898781}
\definecolor{warnbg}{HTML}{FFF6E5}
\definecolor{warnline}{HTML}{EDA100}

\titleformat{\section}{\large\bfseries\color{black}}{\thesection}{0.6em}{}
\titleformat{\subsection}{\normalsize\bfseries}{\thesubsection}{0.6em}{}
\titlespacing*{\section}{0pt}{14pt}{4pt}
\titlespacing*{\subsection}{0pt}{10pt}{2pt}

\pagestyle{fancy}
\fancyhf{}
\renewcommand{\headrulewidth}{0pt}
\newcommand{\shorttitle}{}
\fancyfoot[L]{\small\color{trline}Anchoring in LLMs \textperiodcentered{} \shorttitle}
\fancyfoot[R]{\small\color{trline}\thepage\ / \pageref*{LastPage}}

% figures made by reports/figures.py, and pages of the longer reports' figure files
\newcommand{\rfig}[2][\textwidth]{\includegraphics[width=#1]{reports/figures/#2.pdf}}
\newcommand{\ffig}[2][\textwidth]{\includegraphics[page=#2,width=#1]{full_report_figures.pdf}}
\newcommand{\sfig}[2][\textwidth]{\includegraphics[page=#2,width=#1]{study2_figures.pdf}}
\newcommand{\efig}[2][\textwidth]{\includegraphics[page=#2,width=#1]{events_2026_figures.pdf}}

% title block: \reporttitle{label}{title}{subtitle}{date}
\newcommand{\reporttitle}[4]{%
  \thispagestyle{fancy}%
  {\small\color{accent}\textsf{\textbf{\MakeUppercase{#1}}}}\par\smallskip
  {\LARGE\bfseries #2\par}\smallskip
  {\large #3\par}\medskip
  {\small Hannan Seyfi \textperiodcentered{} Supervisors: Dr.\ Shiva Kamkar and Dr.\ Babak Majidi \textperiodcentered{} Khatam University \textperiodcentered{} #4\par}
  {\small Project: \emph{Adjusting or Retrieving? Testing the Mechanisms of Anchoring in Large Language Models}\par}
  \medskip{\color{trline}\hrule}\medskip}

% plain-language abstract
\newenvironment{plainabstract}{\begin{tcolorbox}[colback=boxblue, colframe=accent, boxrule=0pt, leftrule=3pt, arc=0pt,
  left=9pt, right=9pt, top=6pt, bottom=6pt, title={\textbf{Abstract, in plain words}}, coltitle=black,
  colbacktitle=boxblue, attach title to upper={\par\smallskip}, fonttitle=\normalsize]}{\end{tcolorbox}}

% key numbers: a two-column table
\newenvironment{keynumbers}{\begin{tcolorbox}[colback=white, colframe=trline, boxrule=0.4pt, arc=0pt, left=8pt,
  right=8pt, top=4pt, bottom=4pt, title={\textbf{At a glance}}, coltitle=black, colbacktitle=trgrey,
  fonttitle=\small]\small\begin{tabular}{@{}P{3.6cm}>{\raggedright\arraybackslash}p{\dimexpr\linewidth-3.6cm-2\tabcolsep\relax}@{}}}{\end{tabular}\end{tcolorbox}}

% a key message
\newtcolorbox{keybox}[1][]{breakable, colback=boxblue, colframe=accent, boxrule=0pt, leftrule=3pt,
  arc=0pt, left=8pt, right=8pt, top=4pt, bottom=4pt, fonttitle=\bfseries, coltitle=black,
  title={#1}, attach title to upper={\ }, before skip=8pt, after skip=8pt}
% a caution
\newtcolorbox{cautionbox}[1][]{breakable, colback=warnbg, colframe=warnline, boxrule=0pt, leftrule=3pt,
  arc=0pt, left=8pt, right=8pt, top=4pt, bottom=4pt, fonttitle=\bfseries, coltitle=black,
  title={#1}, attach title to upper={\ }, before skip=8pt, after skip=8pt}
% literal model text
\newtcolorbox{transcript}[1][]{breakable, colback=trgrey, colframe=trline, boxrule=0pt, leftrule=2pt,
  arc=0pt, left=8pt, right=8pt, top=4pt, bottom=4pt, fontupper=\small, fonttitle=\small\bfseries,
  coltitle=black, title={#1}, attach title to upper={\par\smallskip}, before skip=8pt, after skip=8pt}
\newcommand{\who}[1]{\par\smallskip\noindent\textsf{\textbf{\footnotesize\MakeUppercase{#1}}}\quad}
\newcommand{\cut}{\textup{[\dots]}}

% (Overleaf compiles from the project root, so paths are relative to it.)
\usepackage[utf8]{inputenc}
\usepackage[T1]{fontenc}
\usepackage{lmodern}
\usepackage[margin=2.4cm]{geometry}
\usepackage{microtype}
\usepackage{amsmath,amssymb}
\usepackage{booktabs}
\usepackage{tabularx}
\usepackage{array}
\usepackage{enumitem}
\usepackage{graphicx}
\usepackage[table]{xcolor}
\usepackage{textcomp}
\usepackage{csquotes}
\usepackage{newunicodechar}
\newunicodechar{θ}{\ensuremath{\theta}}
\newunicodechar{φ}{\ensuremath{\varphi}}
\newunicodechar{λ}{\ensuremath{\lambda}}
\newunicodechar{Δ}{\ensuremath{\Delta}}
\newunicodechar{≈}{\ensuremath{\approx}}
\newunicodechar{°}{\ensuremath{^\circ}}
\newunicodechar{π}{\ensuremath{\pi}}
\newunicodechar{√}{\ensuremath{\surd}}
\newunicodechar{²}{\ensuremath{^2}}
\newunicodechar{³}{\ensuremath{^3}}
\newunicodechar{−}{\ensuremath{-}}
\newunicodechar{⋅}{\ensuremath{\cdot}}
\usepackage[most]{tcolorbox}
\usepackage[section]{placeins}
\usepackage{titlesec}
\usepackage{fancyhdr}
\usepackage{lastpage}
\usepackage[authoryear,round]{natbib}
\usepackage[colorlinks=true,linkcolor=blue!50!black,citecolor=blue!50!black,urlcolor=blue!50!black]{hyperref}

\setlist{itemsep=2pt,topsep=4pt}
\setlength{\parindent}{0pt}
\setlength{\parskip}{5pt}
\newcolumntype{L}{>{\raggedright\arraybackslash}X}
\newcolumntype{P}[1]{>{\raggedright\arraybackslash}p{#1}}

\definecolor{accent}{HTML}{2A78D6}
\definecolor{boxblue}{HTML}{EEF3FB}
\definecolor{trgrey}{HTML}{F4F4F1}
\definecolor{trline}{HTML}{898781}
\definecolor{warnbg}{HTML}{FFF6E5}
\definecolor{warnline}{HTML}{EDA100}

\titleformat{\section}{\large\bfseries\color{black}}{\thesection}{0.6em}{}
\titleformat{\subsection}{\normalsize\bfseries}{\thesubsection}{0.6em}{}
\titlespacing*{\section}{0pt}{14pt}{4pt}
\titlespacing*{\subsection}{0pt}{10pt}{2pt}

\pagestyle{fancy}
\fancyhf{}
\renewcommand{\headrulewidth}{0pt}
\newcommand{\shorttitle}{}
\fancyfoot[L]{\small\color{trline}Anchoring in LLMs \textperiodcentered{} \shorttitle}
\fancyfoot[R]{\small\color{trline}\thepage\ / \pageref*{LastPage}}

% figures made by reports/figures.py, and pages of the longer reports' figure files
\newcommand{\rfig}[2][\textwidth]{\includegraphics[width=#1]{reports/figures/#2.pdf}}
\newcommand{\ffig}[2][\textwidth]{\includegraphics[page=#2,width=#1]{full_report_figures.pdf}}
\newcommand{\sfig}[2][\textwidth]{\includegraphics[page=#2,width=#1]{study2_figures.pdf}}
\newcommand{\efig}[2][\textwidth]{\includegraphics[page=#2,width=#1]{events_2026_figures.pdf}}

% title block: \reporttitle{label}{title}{subtitle}{date}
\newcommand{\reporttitle}[4]{%
  \thispagestyle{fancy}%
  {\small\color{accent}\textsf{\textbf{\MakeUppercase{#1}}}}\par\smallskip
  {\LARGE\bfseries #2\par}\smallskip
  {\large #3\par}\medskip
  {\small Hannan Seyfi \textperiodcentered{} Supervisors: Dr.\ Shiva Kamkar and Dr.\ Babak Majidi \textperiodcentered{} Khatam University \textperiodcentered{} #4\par}
  {\small Project: \emph{Adjusting or Retrieving? Testing the Mechanisms of Anchoring in Large Language Models}\par}
  \medskip{\color{trline}\hrule}\medskip}

% plain-language abstract
\newenvironment{plainabstract}{\begin{tcolorbox}[colback=boxblue, colframe=accent, boxrule=0pt, leftrule=3pt, arc=0pt,
  left=9pt, right=9pt, top=6pt, bottom=6pt, title={\textbf{Abstract, in plain words}}, coltitle=black,
  colbacktitle=boxblue, attach title to upper={\par\smallskip}, fonttitle=\normalsize]}{\end{tcolorbox}}

% key numbers: a two-column table
\newenvironment{keynumbers}{\begin{tcolorbox}[colback=white, colframe=trline, boxrule=0.4pt, arc=0pt, left=8pt,
  right=8pt, top=4pt, bottom=4pt, title={\textbf{At a glance}}, coltitle=black, colbacktitle=trgrey,
  fonttitle=\small]\small\begin{tabular}{@{}P{3.6cm}>{\raggedright\arraybackslash}p{\dimexpr\linewidth-3.6cm-2\tabcolsep\relax}@{}}}{\end{tabular}\end{tcolorbox}}

% a key message
\newtcolorbox{keybox}[1][]{breakable, colback=boxblue, colframe=accent, boxrule=0pt, leftrule=3pt,
  arc=0pt, left=8pt, right=8pt, top=4pt, bottom=4pt, fonttitle=\bfseries, coltitle=black,
  title={#1}, attach title to upper={\ }, before skip=8pt, after skip=8pt}
% a caution
\newtcolorbox{cautionbox}[1][]{breakable, colback=warnbg, colframe=warnline, boxrule=0pt, leftrule=3pt,
  arc=0pt, left=8pt, right=8pt, top=4pt, bottom=4pt, fonttitle=\bfseries, coltitle=black,
  title={#1}, attach title to upper={\ }, before skip=8pt, after skip=8pt}
% literal model text
\newtcolorbox{transcript}[1][]{breakable, colback=trgrey, colframe=trline, boxrule=0pt, leftrule=2pt,
  arc=0pt, left=8pt, right=8pt, top=4pt, bottom=4pt, fontupper=\small, fonttitle=\small\bfseries,
  coltitle=black, title={#1}, attach title to upper={\par\smallskip}, before skip=8pt, after skip=8pt}
\newcommand{\who}[1]{\par\smallskip\noindent\textsf{\textbf{\footnotesize\MakeUppercase{#1}}}\quad}
\newcommand{\cut}{\textup{[\dots]}}

% (Overleaf compiles from the project root, so paths are relative to it.)
\usepackage[utf8]{inputenc}
\usepackage[T1]{fontenc}
\usepackage{lmodern}
\usepackage[margin=2.4cm]{geometry}
\usepackage{microtype}
\usepackage{amsmath,amssymb}
\usepackage{booktabs}
\usepackage{tabularx}
\usepackage{array}
\usepackage{enumitem}
\usepackage{graphicx}
\usepackage[table]{xcolor}
\usepackage{textcomp}
\usepackage{csquotes}
\usepackage{newunicodechar}
\newunicodechar{θ}{\ensuremath{\theta}}
\newunicodechar{φ}{\ensuremath{\varphi}}
\newunicodechar{λ}{\ensuremath{\lambda}}
\newunicodechar{Δ}{\ensuremath{\Delta}}
\newunicodechar{≈}{\ensuremath{\approx}}
\newunicodechar{°}{\ensuremath{^\circ}}
\newunicodechar{π}{\ensuremath{\pi}}
\newunicodechar{√}{\ensuremath{\surd}}
\newunicodechar{²}{\ensuremath{^2}}
\newunicodechar{³}{\ensuremath{^3}}
\newunicodechar{−}{\ensuremath{-}}
\newunicodechar{⋅}{\ensuremath{\cdot}}
\usepackage[most]{tcolorbox}
\usepackage[section]{placeins}
\usepackage{titlesec}
\usepackage{fancyhdr}
\usepackage{lastpage}
\usepackage[authoryear,round]{natbib}
\usepackage[colorlinks=true,linkcolor=blue!50!black,citecolor=blue!50!black,urlcolor=blue!50!black]{hyperref}

\setlist{itemsep=2pt,topsep=4pt}
\setlength{\parindent}{0pt}
\setlength{\parskip}{5pt}
\newcolumntype{L}{>{\raggedright\arraybackslash}X}
\newcolumntype{P}[1]{>{\raggedright\arraybackslash}p{#1}}

\definecolor{accent}{HTML}{2A78D6}
\definecolor{boxblue}{HTML}{EEF3FB}
\definecolor{trgrey}{HTML}{F4F4F1}
\definecolor{trline}{HTML}{898781}
\definecolor{warnbg}{HTML}{FFF6E5}
\definecolor{warnline}{HTML}{EDA100}

\titleformat{\section}{\large\bfseries\color{black}}{\thesection}{0.6em}{}
\titleformat{\subsection}{\normalsize\bfseries}{\thesubsection}{0.6em}{}
\titlespacing*{\section}{0pt}{14pt}{4pt}
\titlespacing*{\subsection}{0pt}{10pt}{2pt}

\pagestyle{fancy}
\fancyhf{}
\renewcommand{\headrulewidth}{0pt}
\newcommand{\shorttitle}{}
\fancyfoot[L]{\small\color{trline}Anchoring in LLMs \textperiodcentered{} \shorttitle}
\fancyfoot[R]{\small\color{trline}\thepage\ / \pageref*{LastPage}}

% figures made by reports/figures.py, and pages of the longer reports' figure files
\newcommand{\rfig}[2][\textwidth]{\includegraphics[width=#1]{reports/figures/#2.pdf}}
\newcommand{\ffig}[2][\textwidth]{\includegraphics[page=#2,width=#1]{full_report_figures.pdf}}
\newcommand{\sfig}[2][\textwidth]{\includegraphics[page=#2,width=#1]{study2_figures.pdf}}
\newcommand{\efig}[2][\textwidth]{\includegraphics[page=#2,width=#1]{events_2026_figures.pdf}}

% title block: \reporttitle{label}{title}{subtitle}{date}
\newcommand{\reporttitle}[4]{%
  \thispagestyle{fancy}%
  {\small\color{accent}\textsf{\textbf{\MakeUppercase{#1}}}}\par\smallskip
  {\LARGE\bfseries #2\par}\smallskip
  {\large #3\par}\medskip
  {\small Hannan Seyfi \textperiodcentered{} Supervisors: Dr.\ Shiva Kamkar and Dr.\ Babak Majidi \textperiodcentered{} Khatam University \textperiodcentered{} #4\par}
  {\small Project: \emph{Adjusting or Retrieving? Testing the Mechanisms of Anchoring in Large Language Models}\par}
  \medskip{\color{trline}\hrule}\medskip}

% plain-language abstract
\newenvironment{plainabstract}{\begin{tcolorbox}[colback=boxblue, colframe=accent, boxrule=0pt, leftrule=3pt, arc=0pt,
  left=9pt, right=9pt, top=6pt, bottom=6pt, title={\textbf{Abstract, in plain words}}, coltitle=black,
  colbacktitle=boxblue, attach title to upper={\par\smallskip}, fonttitle=\normalsize]}{\end{tcolorbox}}

% key numbers: a two-column table
\newenvironment{keynumbers}{\begin{tcolorbox}[colback=white, colframe=trline, boxrule=0.4pt, arc=0pt, left=8pt,
  right=8pt, top=4pt, bottom=4pt, title={\textbf{At a glance}}, coltitle=black, colbacktitle=trgrey,
  fonttitle=\small]\small\begin{tabular}{@{}P{3.6cm}>{\raggedright\arraybackslash}p{\dimexpr\linewidth-3.6cm-2\tabcolsep\relax}@{}}}{\end{tabular}\end{tcolorbox}}

% a key message
\newtcolorbox{keybox}[1][]{breakable, colback=boxblue, colframe=accent, boxrule=0pt, leftrule=3pt,
  arc=0pt, left=8pt, right=8pt, top=4pt, bottom=4pt, fonttitle=\bfseries, coltitle=black,
  title={#1}, attach title to upper={\ }, before skip=8pt, after skip=8pt}
% a caution
\newtcolorbox{cautionbox}[1][]{breakable, colback=warnbg, colframe=warnline, boxrule=0pt, leftrule=3pt,
  arc=0pt, left=8pt, right=8pt, top=4pt, bottom=4pt, fonttitle=\bfseries, coltitle=black,
  title={#1}, attach title to upper={\ }, before skip=8pt, after skip=8pt}
% literal model text
\newtcolorbox{transcript}[1][]{breakable, colback=trgrey, colframe=trline, boxrule=0pt, leftrule=2pt,
  arc=0pt, left=8pt, right=8pt, top=4pt, bottom=4pt, fontupper=\small, fonttitle=\small\bfseries,
  coltitle=black, title={#1}, attach title to upper={\par\smallskip}, before skip=8pt, after skip=8pt}
\newcommand{\who}[1]{\par\smallskip\noindent\textsf{\textbf{\footnotesize\MakeUppercase{#1}}}\quad}
\newcommand{\cut}{\textup{[\dots]}}

% (Overleaf compiles from the project root, so paths are relative to it.)
\usepackage[utf8]{inputenc}
\usepackage[T1]{fontenc}
\usepackage{lmodern}
\usepackage[margin=2.4cm]{geometry}
\usepackage{microtype}
\usepackage{amsmath,amssymb}
\usepackage{booktabs}
\usepackage{tabularx}
\usepackage{array}
\usepackage{enumitem}
\usepackage{graphicx}
\usepackage[table]{xcolor}
\usepackage{textcomp}
\usepackage{csquotes}
\usepackage{newunicodechar}
\newunicodechar{θ}{\ensuremath{\theta}}
\newunicodechar{φ}{\ensuremath{\varphi}}
\newunicodechar{λ}{\ensuremath{\lambda}}
\newunicodechar{Δ}{\ensuremath{\Delta}}
\newunicodechar{≈}{\ensuremath{\approx}}
\newunicodechar{°}{\ensuremath{^\circ}}
\newunicodechar{π}{\ensuremath{\pi}}
\newunicodechar{√}{\ensuremath{\surd}}
\newunicodechar{²}{\ensuremath{^2}}
\newunicodechar{³}{\ensuremath{^3}}
\newunicodechar{−}{\ensuremath{-}}
\newunicodechar{⋅}{\ensuremath{\cdot}}
\usepackage[most]{tcolorbox}
\usepackage[section]{placeins}
\usepackage{titlesec}
\usepackage{fancyhdr}
\usepackage{lastpage}
\usepackage[authoryear,round]{natbib}
\usepackage[colorlinks=true,linkcolor=blue!50!black,citecolor=blue!50!black,urlcolor=blue!50!black]{hyperref}

\setlist{itemsep=2pt,topsep=4pt}
\setlength{\parindent}{0pt}
\setlength{\parskip}{5pt}
\newcolumntype{L}{>{\raggedright\arraybackslash}X}
\newcolumntype{P}[1]{>{\raggedright\arraybackslash}p{#1}}

\definecolor{accent}{HTML}{2A78D6}
\definecolor{boxblue}{HTML}{EEF3FB}
\definecolor{trgrey}{HTML}{F4F4F1}
\definecolor{trline}{HTML}{898781}
\definecolor{warnbg}{HTML}{FFF6E5}
\definecolor{warnline}{HTML}{EDA100}

\titleformat{\section}{\large\bfseries\color{black}}{\thesection}{0.6em}{}
\titleformat{\subsection}{\normalsize\bfseries}{\thesubsection}{0.6em}{}
\titlespacing*{\section}{0pt}{14pt}{4pt}
\titlespacing*{\subsection}{0pt}{10pt}{2pt}

\pagestyle{fancy}
\fancyhf{}
\renewcommand{\headrulewidth}{0pt}
\newcommand{\shorttitle}{}
\fancyfoot[L]{\small\color{trline}Anchoring in LLMs \textperiodcentered{} \shorttitle}
\fancyfoot[R]{\small\color{trline}\thepage\ / \pageref*{LastPage}}

% figures made by reports/figures.py, and pages of the longer reports' figure files
\newcommand{\rfig}[2][\textwidth]{\includegraphics[width=#1]{reports/figures/#2.pdf}}
\newcommand{\ffig}[2][\textwidth]{\includegraphics[page=#2,width=#1]{full_report_figures.pdf}}
\newcommand{\sfig}[2][\textwidth]{\includegraphics[page=#2,width=#1]{study2_figures.pdf}}
\newcommand{\efig}[2][\textwidth]{\includegraphics[page=#2,width=#1]{events_2026_figures.pdf}}

% title block: \reporttitle{label}{title}{subtitle}{date}
\newcommand{\reporttitle}[4]{%
  \thispagestyle{fancy}%
  {\small\color{accent}\textsf{\textbf{\MakeUppercase{#1}}}}\par\smallskip
  {\LARGE\bfseries #2\par}\smallskip
  {\large #3\par}\medskip
  {\small Hannan Seyfi \textperiodcentered{} Supervisors: Dr.\ Shiva Kamkar and Dr.\ Babak Majidi \textperiodcentered{} Khatam University \textperiodcentered{} #4\par}
  {\small Project: \emph{Adjusting or Retrieving? Testing the Mechanisms of Anchoring in Large Language Models}\par}
  \medskip{\color{trline}\hrule}\medskip}

% plain-language abstract
\newenvironment{plainabstract}{\begin{tcolorbox}[colback=boxblue, colframe=accent, boxrule=0pt, leftrule=3pt, arc=0pt,
  left=9pt, right=9pt, top=6pt, bottom=6pt, title={\textbf{Abstract, in plain words}}, coltitle=black,
  colbacktitle=boxblue, attach title to upper={\par\smallskip}, fonttitle=\normalsize]}{\end{tcolorbox}}

% key numbers: a two-column table
\newenvironment{keynumbers}{\begin{tcolorbox}[colback=white, colframe=trline, boxrule=0.4pt, arc=0pt, left=8pt,
  right=8pt, top=4pt, bottom=4pt, title={\textbf{At a glance}}, coltitle=black, colbacktitle=trgrey,
  fonttitle=\small]\small\begin{tabular}{@{}P{3.6cm}>{\raggedright\arraybackslash}p{\dimexpr\linewidth-3.6cm-2\tabcolsep\relax}@{}}}{\end{tabular}\end{tcolorbox}}

% a key message
\newtcolorbox{keybox}[1][]{breakable, colback=boxblue, colframe=accent, boxrule=0pt, leftrule=3pt,
  arc=0pt, left=8pt, right=8pt, top=4pt, bottom=4pt, fonttitle=\bfseries, coltitle=black,
  title={#1}, attach title to upper={\ }, before skip=8pt, after skip=8pt}
% a caution
\newtcolorbox{cautionbox}[1][]{breakable, colback=warnbg, colframe=warnline, boxrule=0pt, leftrule=3pt,
  arc=0pt, left=8pt, right=8pt, top=4pt, bottom=4pt, fonttitle=\bfseries, coltitle=black,
  title={#1}, attach title to upper={\ }, before skip=8pt, after skip=8pt}
% literal model text
\newtcolorbox{transcript}[1][]{breakable, colback=trgrey, colframe=trline, boxrule=0pt, leftrule=2pt,
  arc=0pt, left=8pt, right=8pt, top=4pt, bottom=4pt, fontupper=\small, fonttitle=\small\bfseries,
  coltitle=black, title={#1}, attach title to upper={\par\smallskip}, before skip=8pt, after skip=8pt}
\newcommand{\who}[1]{\par\smallskip\noindent\textsf{\textbf{\footnotesize\MakeUppercase{#1}}}\quad}
\newcommand{\cut}{\textup{[\dots]}}
